\clearpage
\VerdierPhysicalPageAnchor{0180}%
\begingroup\small\setlength{\parskip}{0pt}\setlength{\emergencystretch}{1.5em}
\setlength{\abovedisplayskip}{4pt plus 2pt minus 2pt}%
\setlength{\belowdisplayskip}{4pt plus 2pt minus 2pt}%

\VerdierSemanticAnchor{ch3d-1-2-8-duality-cont}
此后，除非另有说明，我们将借助同构 $i$ 把范畴 $D(\mathcal A)^o$ 与
$D(\mathcal A^o)$ 视为同一个范畴。注意，这一等同把范畴
$D^b(\mathcal A)^o$ 与范畴 $D^b(\mathcal A^o)$ 等同，把范畴
$D^+(\mathcal A)^o$ 与范畴 $D^-(\mathcal A^o)$ 等同，并把范畴
$D^-(\mathcal A)^o$ 与范畴 $D^+(\mathcal A^o)$ 等同。

\VerdierSemanticAnchor{ch3d-1-2-9}
\medskip\noindent\textbf{1.2.9.}\quad
典范函子 $\mathcal A\longrightarrow K(\mathcal A)$（第二章，1.3.3）与函子
$Q\colon K(\mathcal A)\longrightarrow D(\mathcal A)$ 的复合给出一个充分忠实
函子 $\mathcal A\longrightarrow D(\mathcal A)$。事实上，设 $X$ 和 $Y$ 是
$\mathcal A$ 的两个对象。$D(\mathcal A)$ 中从 $X$ 到 $Y$ 的一个态射是
（第二章，2.2.2）$K(\mathcal A)$ 中一个图式的类 $\overline D$：

\VerdierSemanticAnchor{ch3d-1-2-9-roof}
\[
\overline D=X\xleftarrow{\ s\ }Z^\bullet\xrightarrow{\ m\ }Y,
\]
其中 $s$ 是 $K(\mathcal A)$ 中的一个拟同构。因此，
$s\colon Z^\bullet\longrightarrow X$ 诱导同构
$H^0(s)\colon H^0(Z^\bullet)\xrightarrow{\sim}H(X)=X$，并且立即验证
$\overline D=Q\bigl(H^0(m)H^0(s)^{-1}\bigr)$。这表明同态
\[
\operatorname{Hom}_{\mathcal A}(X,Y)\longrightarrow
\operatorname{Hom}_{D(\mathcal A)}(X,Y)
\]
是满射。下面证明它是单射。$\mathcal A$ 中的态射
$f\colon X\longrightarrow Y$ 在范畴 $D(\mathcal A)$ 中定义零态射，当且仅当
存在 $K(\mathcal A)$ 中的拟同构 $s\colon Z^\bullet\longrightarrow X$，使得
$fs=0$。由此得到 $H^0(f)H^0(s)=0$。由于 $H^0(s)$ 是同构，遂有
$H^0(f)=f=0$。

还要注意，函子 $\mathcal A\longrightarrow D(\mathcal A)$ 在对象集上也是
单射，因而将 $\mathcal A$ 实现为 $D(\mathcal A)$ 的满子范畴。此后，始终
通过这一函子把范畴 $\mathcal A$ 看作其导出范畴的一个子范畴。典范上同调函子
$D(\mathcal A)\xrightarrow{H}\mathcal A$ 限制在子范畴 $\mathcal A$ 上就是
恒等函子。

\VerdierSemanticAnchor{ch3d-1-2-10}
\medskip\noindent\textbf{命题 1.2.10.}\quad
\textit{a) $D(\mathcal A)$ 的对象 $X$ 同构于 $\mathcal A$ 的一个对象，
当且仅当对每个 $i\neq0$ 都有 $H^i(X)=0$。此时它典范同构于 $H^0(X)$。}

\noindent\textit{b) 设 $X$ 和 $Y$ 是 $\mathcal A$ 的两个对象。则有：}
\[
\operatorname{Ext}_{\mathcal A}^{n}(X,Y)=0\quad\text{当 }n<0,
\qquad
\operatorname{Ext}_{\mathcal A}^{0}(X,Y)
=\operatorname{Hom}_{\mathcal A}(X,Y).
\]

\VerdierSemanticAnchor{ch3d-1-2-10-proof}
证明 a)。当 $X$ 是 $\mathcal A$ 的一个对象时，对 $i\neq0$ 有
$H^i(X)=0$ 是显然的。设 $X$ 是 $D(\mathcal A)$ 的一个对象，并且当
$i\neq0$ 时，$H^i(X)$ 为零。
\endgroup
